%-----------------------------------------------------
%
%    Content
%
%-----------------------------------------------------
\sectionfix
\section*{Abstract}
Recent advances in the HCI community have proven that on-body sensors changes the way we interact with our environment. By using our body, we obtain a better understanding on the correlation between perception and manipulation of the digital world. An initial exploration in this direction is the development of design tools for sensors. It helps maintaining the cost of fabrication and improves design iterations for better performance.\\
In this thesis, we propose a novel system (refered to as Scansify) for personalized design and rapid prototyping of thin-layered sensors. Using a low budget RBG-D camera, we were able to extract the finer details of the human body and enable the user to make annotations on top of its 3D model.\\
We provide design principles and detailed workflow description of the graphical user interface. Additionally, a user study is conducted to prove the feasibility and usability of the tool. Feedback in a semi-unstructured interview has been evaluated to obtain critical data for future iterations. Thereby, results illustrate that Scansify is a valuable tool for customized design processes of on-skin sensors.
